\documentclass[10pt,a4paper]{article}

\usepackage[utf8]{inputenc}
\usepackage[T1]{fontenc}
\usepackage{lmodern}
\usepackage{microtype}
\usepackage[margin=2.5cm]{geometry}
\usepackage{graphicx}
\usepackage{booktabs}
\usepackage{multirow}
\usepackage{amsmath, amssymb}
\usepackage[round]{natbib}
\usepackage{hyperref}
\usepackage{cleveref}
\usepackage{xcolor}
\usepackage{caption}

\hypersetup{colorlinks=true, linkcolor=black, citecolor=black, urlcolor=blue}

\title{Benchmarking ColVision retrieval models within the \textsc{mmore} RAG pipeline:\\
a scaling and multilingual study on medical documents}
\author{Mathieu Bonnet \\ \texttt{mathieu.bonnet@telecom-sudparis.eu}}
\date{\today}

\begin{document}
\maketitle

\begin{abstract}
We present an open-source benchmark suite that compares five ColVision
multi-vector retrieval models (ColPali, ColQwen2, ColQwen2.5, ColQwen3 and a
ColGemma3 variant) within a single, fixed RAG pipeline (\textsc{mmore} PR
\#305). Unlike per-model evaluations on \textsc{ViDoRe}, we hold the index,
the reranker, the embedding loader and the retrieval contract constant, and
vary only the visual encoder. Two studies are reported: (A)~an
\emph{English-only scaling study} on a figure-rich PubMed Central Open Access
subset at $\{100, 1k, 10k, 50k\}$ pages, and (B)~a \emph{multilingual study}
on six languages (English, French, German, Spanish, Simplified Chinese,
Arabic) at a fixed corpus size, including one logographic and one RTL
script. All cells are run with three random seeds and reported with 95\%
bootstrap confidence intervals; pairwise comparisons use Wilcoxon signed-rank
with Holm correction. Queries are generated by an inverse prompt that asks
\textsc{Meditron}-70B to author the questions a page would uniquely answer; we
validate this methodology by correlating its yield against an annotated
healthcare subset of \textsc{ViDoRe}. We release the benchmark code, the
corpus manifest with per-PDF SHA-256, and the synthetic query sets.
\end{abstract}

\section{Introduction}

Vision-language retrieval over document collections has shifted, in less than
two years, from a research curiosity to a production primitive. The ColVision
line of work --- starting with ColPali~\cite{faysse2024colpali} and continued
by ColQwen2, ColQwen2.5, ColQwen3 and ColGemma3 variants --- has demonstrated
that multi-vector page-level embeddings, combined with late-interaction
reranking inherited from ColBERT~\cite{khattab2020colbert}, can outperform
text-only OCR$\rightarrow$BM25 pipelines on figure-rich corpora. Yet the
public evaluation surface for these encoders is fragmented. Each model is
released with its own \textsc{ViDoRe}~\cite{faysse2024vidore} scores, computed
on the authors' own software stack, often against subtly different versions
of the benchmark. \textsc{MMEB}~\cite{wei2025mmeb} reduces some of this
variance for general multimodal embedding tasks, but does not target document
retrieval and assumes a flat single-vector contract incompatible with
late-interaction encoders.

The result is that practitioners considering ColVision for a real-world
retrieval-augmented-generation (RAG) deployment cannot read off, from
published numbers, which encoder best suits their corpus and constraints.
Three questions are answered nowhere in the public literature:

\begin{enumerate}
  \item \textbf{Cross-model parity.} How do the five ColVision models compare
    when the index, the reranker, the embedding loader and the retrieval
    contract are all held strictly constant?
  \item \textbf{Scaling.} How does each model's retrieval quality, latency
    and memory degrade as the index grows from a few hundred to tens of
    thousands of pages --- the regime an actual hospital or pharmaceutical
    knowledge base operates in?
  \item \textbf{Multilingual robustness.} How much do the models generalise
    across scripts unrelated to their pre-training distribution, including
    logographic (Simplified Chinese) and RTL connected scripts (Arabic)?
\end{enumerate}

We address all three by building a fixed-pipeline benchmark on top of the
\textsc{mmore} project~\cite{mmore2024} (PR \#305), which integrates the five
encoders behind a single CLI. We invoke the \textsc{mmore} CLI as a
subprocess rather than re-implementing its pipeline: every measurement we
report reflects the exact code path a real \textsc{mmore} user would
exercise. The benchmark code, configuration manifests with per-PDF SHA-256,
synthetic query sets and result records (in an
\textsc{MTEB}-compatible~\cite{muennighoff2023mteb} JSON schema) are
released alongside this report.

\section{Related work}

\paragraph{Visual document retrieval.} ColPali~\cite{faysse2024colpali}
introduced page-level multi-vector retrieval over rendered PDF images, using
PaliGemma as the visual encoder and a ColBERT-style late-interaction scorer.
Subsequent ColQwen2/2.5/3 variants swap the encoder for the Qwen-VL
series~\cite{bai2023qwenvl}, generally trading throughput for quality.
ColGemma3 variants (e.g.~\textit{ColNetraEmbed}) target multilingual
performance through the Gemma-3 backbone~\cite{team2024gemma3}. All four
families are evaluated by their authors on
\textsc{ViDoRe}~\cite{faysse2024vidore} and its v2
extension~\cite{faysse2025vidorev2}, but typically with different software
stacks, different reranker implementations, and different batch-size /
precision configurations, making cross-paper comparison unreliable.

\paragraph{Multimodal embedding benchmarks.}
\textsc{MMEB}~\cite{wei2025mmeb} standardises the multimodal embedding
evaluation surface across 36 datasets, but operates on single-vector
embeddings and on tasks (classification, VQA, grounding) that do not require
page-level multi-vector retrieval. \textsc{MTEB}~\cite{muennighoff2023mteb}
plays the same role for text and provides the result-format conventions we
adopt for this work; ColVision models fall outside its scope.

\paragraph{RAG generation metrics.}
\textsc{RAGAS}~\cite{es2023ragas} proposes LLM-as-judge metrics for
retrieval-augmented generation --- faithfulness, answer relevancy, context
precision and context recall --- which we adopt unchanged. We deliberately
use a single judge (Meditron-70B~\cite{chen2023meditron}) and call out the
inter-judge-agreement gap as a limitation rather than averaging across judges.

\paragraph{Synthetic query generation.} The idea of asking a strong LLM to
author the queries a page would uniquely answer dates back to the
\textsc{InPars}~\cite{bonifacio2022inpars} and \textsc{Promptagator}
~\cite{dai2022promptagator} lines of work in text retrieval. We follow the
same principle in the document-image setting, prompt Meditron-70B with the
page's extracted text plus visual hints, and add a second-pass ambiguity
filter modelled on the same judge. The methodology is validated (Phase 4)
against an annotated subset of \textsc{ViDoRe} healthcare.

\section{Methodology}

\subsection{Pipeline under test}

\textsc{mmore} PR \#305 wires the ColVision encoders into a single ingest /
retrieve pipeline. Each PDF page is rendered to a 200-dpi PNG via
\textsc{PyMuPDF}~\cite{pymupdf}, then encoded by the chosen ColVision model
into a multi-vector (token-level) embedding tensor. Embeddings are
serialised to Parquet and inserted into a Milvus~\cite{wang2021milvus}
collection (\textsc{FLAT} index, inner-product metric) keyed by the page
identifier. Retrieval issues an embedding-only top-$k$ over Milvus, then
re-ranks the candidates with a ColBERT-style MaxSim score~\cite{khattab2020colbert}
computed in a CPU thread pool. The resulting ranked list is returned to the
caller through a LangChain \textit{Retriever} contract; the generation step,
when used, feeds the top-$k$ pages to a separate generator LLM.

We hold every component of this pipeline constant across models. The only
variables are (a)~the ColVision encoder, (b)~the corpus, and (c)~the random
seed. The Milvus collection name is derived from the model identifier so
that different embedding dimensions do not collide.

\subsection{Models}

The five encoders under evaluation are listed in Table~\ref{tab:models}.
Embedding dimensions differ across families (128 for the ColPali/ColQwen
line, 320 for ColQwen3, 768 for the ColGemma3 variant), which the pipeline
absorbs by giving each model its own Milvus collection. The batch sizes and
precision are taken from each model's published recommendation; we never
re-tune them.

\begin{table}[h]
  \centering
  \caption{ColVision encoders evaluated in this work.}
  \label{tab:models}
  \begin{tabular}{lllrll}
    \toprule
    Id & Family & HuggingFace name & Dim & Batch & Dtype \\
    \midrule
    \texttt{colpali\_v1\_3}        & ColPali   & \texttt{vidore/colpali-v1.3}        & 128 & 4 & bf16 \\
    \texttt{colqwen2\_v1\_0}       & ColQwen2  & \texttt{vidore/colqwen2-v1.0}       & 128 & 4 & bf16 \\
    \texttt{colqwen2\_5\_v0\_2}    & ColQwen2.5& \texttt{vidore/colqwen2.5-v0.2}     & 128 & 4 & bf16 \\
    \texttt{colqwen3\_v0\_1}       & ColQwen3  & \texttt{vidore/colqwen3-v0.1}       & 320 & 2 & bf16 \\
    \texttt{colgemma3\_colnetra}   & ColGemma3 & \texttt{Cognitive-Lab/ColNetraEmbed}& 768 & 2 & bf16 \\
    \bottomrule
  \end{tabular}
\end{table}

\subsection{Corpora}

\paragraph{Track A (English scaling).} We assemble a figure-rich English
medical corpus from PubMed Central Open Access~\cite{pmcoa}. Articles are
filtered by visual density: for each PDF we compute the fraction of total
page area occupied by figures and tables (\textsc{PyMuPDF} bounding boxes),
and keep PDFs whose density is at least 30\%. The retained PDFs are sampled,
stratified by their MeSH tag, to four palier sizes: 100, 1{,}000, 10{,}000
and 50{,}000 pages. Each palier is a strict subset of the next, so a model
indexing the 50k-page palier has already paid the indexing cost for the
smaller ones --- this is the regime of an organisation that grows its
knowledge base incrementally.

\paragraph{Track B (multilingual).} For each of six languages --- English,
French, German, Spanish, Simplified Chinese, Arabic --- we target
$\approx$1{,}500 pages from a language-appropriate open-access medical
source: PMC-OA for English, HAL+Cairn for French, Thieme OA for German,
SciELO for Spanish, CNKI OA for Chinese, the Saudi Medical Journal for
Arabic. The same visual-density filter is applied. Chinese and Arabic are
selected deliberately to stress the encoders along two orthogonal axes: a
logographic script that shares no characters with the encoders' Latin-heavy
pre-training distribution, and an RTL connected script that breaks the
character-segmentation assumptions baked into most OCR-based pre-training
proxies.

\paragraph{Reproducibility.} Every corpus is shipped as a
\texttt{CorpusManifest} JSON: for each PDF the manifest records its
SHA-256 hash, page count, detected language, visual density, MeSH tags and
source. The \texttt{bcv-corpus verify} command re-hashes the local files to
detect drift; the manifest hash is recorded in every benchmark result so
that any reported number can be traced back to the bytes that produced it.

\subsection{Synthetic query generation}

We generate one query set per (track, language, palier) tuple. For each
page that passes the visual-density filter, we extract the page text with
PyMuPDF, build a prompt that includes the text plus a structured visual hint
(``this page contains $\geq$1 figure / table''), and ask Meditron-70B (via
vLLM, OpenAI-compatible endpoint) to author exactly two clinical questions
that can be answered \emph{only} by consulting that specific page. The
prompt asks the LLM to flag each question with a \texttt{requires\_visual}
boolean; we keep only those flagged \texttt{true}, ensuring the benchmark
exercises the visual dimension rather than collapsing to a text-retrieval
proxy.

A second pass --- the \emph{ambiguity filter} --- re-presents each
candidate question to the same judge along with the page text, and asks for
a single confidence score in $[0, 1]$ representing whether the question is
(a)~precise enough to admit a single best page-source, and (b)~genuinely
tied to that particular page. Queries below a 0.8 confidence threshold are
dropped. The kept queries are serialised as a \texttt{QuerySet} JSONL whose
SHA-256 is recorded in every downstream result. Volumes are sized to
$\sim$2{,}000 retained queries on Track A (spread across paliers) and
$\sim$200 per language on Track B.

\subsection{Methodology validation}

Synthetic queries inherit the biases of the LLM that authored them, in
particular a tendency to over-rely on textual cues that happen to be unique
to a page. We bound this bias by validating the methodology against an
annotated reference: we take a $\sim$50-document healthcare slice of
\textsc{ViDoRe} v2 for which human-authored queries and gold pages exist,
re-generate queries on the same documents with our pipeline, and compute
nDCG@5 on both query sets under an identical retrieval configuration. We
report the Pearson and Spearman correlations between the two per-document
nDCG@5 vectors; the benchmark gates on a Pearson coefficient
$\geq$ 0.85. If the gate fails we iterate on the generation prompt before
running any benchmark cell, on the assumption that a methodology that does
not correlate with human judgement on a known-good subset cannot be trusted
on novel corpora.

\section{Track A --- scaling study}
\subsection{Setup}

For each of the five models and four paliers we run three random seeds. A
cell consists of: (1) processing the palier PDFs with
\texttt{mmore colpali process}, (2) indexing the resulting embeddings with
\texttt{mmore colpali index}, (3) running batched retrieval with
\texttt{mmore colpali retrieve} over the palier's query set. Each step's
duration, GPU memory peak (via \textsc{NVML}) and per-query latency are
captured into a \texttt{BenchmarkRecord}. Sixty cells in total
(5~models~$\times$~4~paliers~$\times$~3~seeds) are submitted as a SLURM job
array. Hardware: a single H100 GPU per cell, 8 CPU cores, 64~GB RAM.

\subsection{Retrieval results}
% Figure: scaling curves with CI bands.
% Table: nDCG@5 [CI] per model x palier.

\subsection{Generation results}

\subsection{Performance results}
% Latency p50/p95/p99, throughput, GPU memory peak vs corpus size.

\subsection{Statistical analysis}
% Pairwise Wilcoxon, Holm correction, effect sizes.

\section{Track B --- multilingual study}
\subsection{Setup}

For each model and language we run three seeds at the fixed
$\sim$1{,}500-page corpus size. The remaining instrumentation is identical
to Track A. Ninety cells (5~$\times$~6~$\times$~3) are submitted as a
single SLURM array. The English cell of Track~B uses a disjoint corpus
sample from Track~A so that the multilingual results do not implicitly
benchmark against memorised content.

\subsection{Per-language results}
% Heatmap, language gap vs EN baseline.

\subsection{Discussion of script effects (CJK, RTL)}

\section{Discussion}

\section{Limitations}
\begin{itemize}
  \item \textbf{Synthetic queries.} Queries are authored by Meditron-70B
    rather than by clinicians. The Phase~4 correlation gate ensures the
    pipeline tracks human judgement on a known-good subset, but a residual
    bias toward LLM-typical question shapes is inevitable. A
    human-annotated subset of $\sim$50 questions per language is a natural
    follow-up.
  \item \textbf{Single LLM judge.} Both query filtering and RAGAS metrics
    use a single Meditron-70B judge. Inter-judge agreement against, e.g.,
    Claude or a second open-weights model is not measured here; this is an
    explicit future-work item, not a quality claim.
  \item \textbf{Corpus skew.} PMC-OA is US-centric and EN-dominant; the six
    Track~B sources differ in domain mix, image quality and license terms.
    We control for visual density but cannot equalise editorial style.
  \item \textbf{Single GPU class.} All cells run on a single H100. We do
    not measure how each encoder behaves under tighter VRAM constraints
    (A100~40GB, T4) or on mixed-precision fallbacks.
  \item \textbf{One pipeline.} The benchmark holds the pipeline constant
    \emph{by design}, but the absolute numbers we report are conditional on
    \textsc{mmore}'s specific implementation choices (Milvus \textsc{FLAT}
    index, CPU MaxSim reranker, top-$k$=10 default). A different RAG stack
    may reorder the models.
\end{itemize}

\section{Conclusion}

\bibliographystyle{plainnat}
\bibliography{references}

\end{document}
